% =============================================================================
% Guía de Ejecución en Windows del script del Datawarehouse (Avance N°1)
% Documento aparte para que el profesor replique la creación del modelo.
% Compilar con: latexmk guia_ejecucion_windows.tex (o lualatex/pdflatex dos
% veces). Es autónomo: si img/logo_utem.png no está, el encabezado sale sin logo.
% Los resultados esperados (filas por consulta) se verificaron en
% PostgreSQL 17.11 el 05-11-2026 (ver docs/bitacora.md).
% =============================================================================
\documentclass[11pt,letterpaper]{article}

\usepackage{iftex}
\ifPDFTeX
  \usepackage[utf8]{inputenc}
  \usepackage[T1]{fontenc}
  \usepackage{lmodern}
\else
  \usepackage{fontspec}
\fi
\usepackage[spanish,es-tabla,es-noshorthands,es-lcroman]{babel}
\usepackage[letterpaper,left=2.5cm,right=2.5cm,bottom=2.5cm,top=3.6cm,
            headheight=44pt,headsep=12pt]{geometry}
\usepackage{microtype}
\usepackage{parskip}
\usepackage{graphicx}
\usepackage{float}
\usepackage{array}
\usepackage{booktabs}
\usepackage{tabularx}
\usepackage{enumitem}
\usepackage{xcolor}
\usepackage{fancyhdr}
\usepackage{listings}
\usepackage{hyperref}

\definecolor{azulutem}{HTML}{004EAA}
\definecolor{verdeutem}{HTML}{78BF26}
% Enlaces en negro, sin recuadros; siguen siendo pulsables en el PDF
\hypersetup{colorlinks=true, linkcolor=black, citecolor=black,
            urlcolor=black, filecolor=black,
            pdftitle={Guía de Ejecución en Windows: Datawarehouse de Transportes Rapiditos S.A.},
            pdfauthor={Grupo 1, Inteligencia de Negocios, UTEM}}
\setlength{\emergencystretch}{2em}

% Nombres de archivos, tablas y comandos cortos en monoespaciado
\newcommand{\bd}[1]{\texttt{\detokenize{#1}}}

% Bloques de comandos para copiar y pegar
\lstdefinestyle{consola}{
  basicstyle=\ttfamily\small, columns=fullflexible, keepspaces=true,
  frame=single, rulecolor=\color{gray!60}, backgroundcolor=\color{gray!6},
  breaklines=true, showstringspaces=false, upquote=true,
  aboveskip=4pt, belowskip=8pt}
\lstset{style=consola}

\newcommand{\institucion}{%
  {\renewcommand{\arraystretch}{1}%
  \begin{tabular}{@{}m{1.0cm}@{\hspace{0.5em}}m{10cm}@{}}
    \IfFileExists{img/logo_utem.png}{\includegraphics[height=1.2cm]{img/logo_utem.png}}{} &
    \small Universidad Tecnológica Metropolitana\newline
    Facultad de Ingeniería\newline
    Escuela de Informática y Computación
  \end{tabular}}}

\pagestyle{fancy}
\fancyhf{}
\fancyhead[L]{\institucion}
\fancyfoot[R]{\thepage}
\renewcommand{\headrulewidth}{0pt}

\begin{document}

\begin{center}
  {\Large\bfseries Guía de Ejecución en Windows\par}
  \medskip
  {\large Script del Datawarehouse de Transportes Rapiditos S.A.\par}
  \medskip
  Inteligencia de Negocios (INFB8101), Sección 412 -- Proyecto Semestral, Avance N°1\par
  Grupo 1: Ignacio Ahumada Poque, Leonardo Solís Ortiz, Vicente Fuentes
  Rodríguez y Gabriel García Vega\par
  Noviembre de 2026
\end{center}

\section{Propósito}

Esta guía permite replicar en Windows la creación del modelo multidimensional
estrella del Avance N°1 y comprobar el resultado. El script \bd{01_ddl.sql}
está escrito para \textbf{PostgreSQL 17} (se probó en la versión 17.11) y
crea, \emph{sin datos}, el esquema \bd{dwh} con la tabla de hechos
\bd{fact_venta} y sus seis dimensiones.

Los dos scripts son \textbf{autónomos}: no leen ni escriben archivos, no
dependen de otros archivos y funcionan desde cualquier carpeta. Basta con
guardar \bd{01_ddl.sql} y \bd{99_validacion.sql} en una misma carpeta
cualquiera del equipo, por ejemplo \emph{Descargas}. En esta guía, esa carpeta
se llama \emph{carpeta de los scripts}.

\begin{table}[H]
  \centering
  \caption{Archivos que se usan en esta guía}
  \small
  \begin{tabularx}{\linewidth}{@{}>{\raggedright\arraybackslash}p{3.4cm}X@{}}
    \toprule
    \textbf{Archivo} & \textbf{Contenido} \\
    \midrule
    \bd{01_ddl.sql} & Crea el esquema \bd{dwh}: 6 dimensiones y la tabla de
      hechos, con sus llaves, reglas \bd{CHECK}, índices y comentarios \\
    \bd{99_validacion.sql} & 9 consultas al catálogo que verifican lo creado \\
    \bottomrule
  \end{tabularx}
\end{table}

Se ofrecen dos opciones. La \textbf{opción A} instala PostgreSQL y pgAdmin
directamente en Windows y es la más simple. La \textbf{opción B} no instala
PostgreSQL: lo ejecuta en un contenedor de Docker con la misma imagen
(\bd{postgres:17.11}) que usó el grupo.

\textbf{Importante:} el script no crea bases de datos ni usuarios. Se ejecuta
dentro de una base existente y, al comenzar, \emph{borra y vuelve a crear} el
esquema \bd{dwh} de esa base. Por eso se recomienda ejecutarlo en una base
nueva y dedicada, como la base \bd{rapiditos} de esta guía.

\section{Opción A: PostgreSQL 17 y pgAdmin instalados en Windows}

\subsection*{A.1. Instalar PostgreSQL 17}

\begin{enumerate}[leftmargin=*]
  \item Descargar el instalador para Windows desde
        \url{https://www.postgresql.org/download/windows/} (enlace
        \emph{Download the installer}, de EDB). Elegir la versión
        \textbf{17.x} para \emph{Windows x86-64}.
  \item Ejecutar el instalador. En \emph{Select Components} dejar marcados
        \emph{PostgreSQL Server}, \emph{pgAdmin 4} y \emph{Command Line Tools}.
        \emph{Stack Builder} no es necesario.
  \item Definir la contraseña del superusuario \bd{postgres} y anotarla.
  \item Mantener el puerto \bd{5432}. Si XAMPP está instalado no hay conflicto,
        porque MySQL/MariaDB usa el puerto 3306.
  \item Aceptar la configuración regional por omisión y terminar la
        instalación. Al final se puede desmarcar la opción de abrir
        Stack Builder.
\end{enumerate}

\subsection*{A.2. Crear la base de datos}

\begin{enumerate}[leftmargin=*]
  \item Abrir \textbf{pgAdmin 4} desde el menú Inicio. En el panel izquierdo,
        expandir \emph{Servers} \textrightarrow{} \emph{PostgreSQL 17} e
        ingresar la contraseña de \bd{postgres}.
  \item Clic derecho en \emph{Databases} \textrightarrow{} \emph{Create}
        \textrightarrow{} \emph{Database\ldots}. En \emph{Database} escribir
        \bd{rapiditos}. En la pestaña \emph{Definition}, verificar que
        \emph{Encoding} sea \bd{UTF8}. Presionar \emph{Save}.
\end{enumerate}

\subsection*{A.3. Ejecutar el script}

\begin{enumerate}[leftmargin=*]
  \item Seleccionar la base \bd{rapiditos} en el panel izquierdo y abrir
        \emph{Tools} \textrightarrow{} \emph{Query Tool}.
  \item Con el ícono de carpeta (\emph{Open File}) de la barra del Query Tool,
        abrir \bd{01_ddl.sql} desde la carpeta de los scripts.
  \item Sin seleccionar texto, presionar \textbf{F5}: se ejecuta el script
        completo. En la pestaña \emph{Messages} debe aparecer \bd{COMMIT} y
        \emph{Query returned successfully}. La primera vez aparece además el
        aviso \bd{NOTICE: schema "dwh" does not exist, skipping} (o su
        equivalente en español), que es normal: el script intenta borrar un
        esquema que todavía no existe.
  \item Clic derecho en la base \bd{rapiditos} \textrightarrow{}
        \emph{Refresh}. En \emph{Schemas} \textrightarrow{} \emph{dwh}
        \textrightarrow{} \emph{Tables} deben aparecer las 7 tablas.
\end{enumerate}

El script se puede ejecutar todas las veces que se quiera: cada vez deja el
mismo resultado. Si una sentencia fallara, la transacción completa se
revierte y no queda nada a medias.

\subsection*{A.4. Verificar el resultado}

\begin{enumerate}[leftmargin=*]
  \item En el Query Tool, abrir \bd{99_validacion.sql} desde la carpeta de los
        scripts.
  \item pgAdmin muestra solo el resultado de la última consulta cuando se
        ejecuta un archivo completo. Para ver cada una, seleccionar con el
        mouse el texto de una consulta (desde su comentario numerado hasta el
        punto y coma) y presionar \textbf{F5}: se ejecuta solo la selección.
  \item Comparar la cantidad de filas de cada consulta con la
        Tabla~\ref{tab:esperado}.
\end{enumerate}

\begin{table}[H]
  \centering
  \caption{Resultado esperado de \bd{99_validacion.sql}}
  \label{tab:esperado}
  \small
  \begin{tabularx}{\linewidth}{@{}cXc@{}}
    \toprule
    \textbf{N°} & \textbf{Consulta} & \textbf{Filas} \\
    \midrule
    1 & Versión del motor (PostgreSQL 17) y codificación de la base (UTF8) & 1 \\
    2 & Tablas del esquema \bd{dwh} con su comentario & 7 \\
    3 & Columnas de cada tabla & 58 \\
    4 & Llaves primarias (7) y valores únicos (7) & 14 \\
    5 & Llaves foráneas, todas en \bd{fact_venta} & 7 \\
    6 & Restricciones \bd{CHECK} & 20 \\
    7 & Índices: 7 de llaves primarias, 7 de valores únicos y 7 sobre llaves foráneas & 21 \\
    8 & Resumen por tabla (columnas y llaves que salen y llegan a cada tabla) & 7 \\
    9 & Restricciones e índices que no siguen la norma de nomenclatura & \textbf{0} \\
    \bottomrule
  \end{tabularx}
\end{table}

\subsection*{A.5. Alternativa por línea de comandos (psql)}

El instalador deja \bd{psql} en
\lstinline[basicstyle=\ttfamily]|C:\Program Files\PostgreSQL\17\bin|. Para
abrir PowerShell en la carpeta de los scripts, abrir esa carpeta en el
Explorador de archivos, escribir \bd{powershell} en la barra de direcciones y
presionar Enter. Luego:

\begin{lstlisting}
chcp 65001
$env:PGCLIENTENCODING = "UTF8"
$psql = "C:\Program Files\PostgreSQL\17\bin\psql.exe"
& $psql -U postgres -d rapiditos -v ON_ERROR_STOP=1 -f 01_ddl.sql
& $psql -U postgres -d rapiditos -P pager=off -f 99_validacion.sql
\end{lstlisting}

\bd{psql} pide la contraseña de \bd{postgres}. Las dos primeras líneas hacen
que la consola muestre bien las tildes; \bd{psql} puede advertir que la página
de códigos de la consola difiere de la de Windows, y ese aviso se puede
ignorar. El script fija por sí mismo la codificación UTF-8 en su primera
línea, por lo que los datos se guardan bien aunque se omitan esas dos líneas.
La opción \bd{-P pager=off} evita que la salida larga se detenga esperando una
tecla. A diferencia de pgAdmin, \bd{psql} muestra el resultado de las 9
consultas de validación seguidas.

\section{Opción B: PostgreSQL en Docker, sin instalarlo}

Requiere \textbf{Docker Desktop} para Windows, instalado según
\url{https://docs.docker.com/desktop/setup/install/windows-install/}, y abierto
antes de continuar. No se necesita ningún otro archivo: los comandos usan solo
los dos scripts.

\subsection*{B.1. Crear el contenedor con la base de datos}

Abrir PowerShell en la carpeta de los scripts, como en el paso A.5, y
ejecutar el comando siguiente, reemplazando \bd{una_clave} por una contraseña a
elección. El acento grave (\lstinline|`|) al final de una línea es la
continuación de línea de PowerShell.

\begin{lstlisting}
docker run --name rapiditos-postgres `
    -e POSTGRES_PASSWORD=una_clave -e POSTGRES_DB=rapiditos `
    -p 127.0.0.1:5432:5432 -d postgres:17.11
\end{lstlisting}

La primera vez se descarga la imagen \bd{postgres:17.11}. La base
\bd{rapiditos} se crea sola. Antes de seguir hay que esperar a que el servidor
termine de iniciarse: repetir el comando siguiente hasta que la última línea
diga \bd{database system is ready to accept connections}.

\begin{lstlisting}
docker logs rapiditos-postgres
\end{lstlisting}

\subsection*{B.2. Ejecutar el script y verificar}

Los scripts se copian al contenedor y se ejecutan con el \bd{psql} que trae la
imagen; dentro del contenedor no se pide contraseña.

% El bloque va en una minipage para que no se parta entre dos páginas
\noindent\begin{minipage}{\linewidth}
\begin{lstlisting}
chcp 65001
docker cp 01_ddl.sql rapiditos-postgres:/tmp/01_ddl.sql
docker cp 99_validacion.sql rapiditos-postgres:/tmp/99_validacion.sql
docker exec rapiditos-postgres psql -U postgres -d rapiditos `
    -v ON_ERROR_STOP=1 -f /tmp/01_ddl.sql
docker exec rapiditos-postgres psql -U postgres -d rapiditos `
    -P pager=off -f /tmp/99_validacion.sql
\end{lstlisting}
\end{minipage}

La ejecución del script debe terminar en \bd{COMMIT}, y la validación debe
coincidir con la Tabla~\ref{tab:esperado}. \bd{chcp 65001} solo sirve para que
la consola muestre bien las tildes.

Si además se quiere revisar la base con pgAdmin (por ejemplo, el de la
opción A), se registra un servidor con host \bd{localhost}, el puerto
publicado (\bd{5432}), usuario \bd{postgres} y la contraseña elegida en el
paso B.1. Al terminar, \bd{docker rm -f rapiditos-postgres} elimina el
contenedor junto con la base.

\section{Problemas frecuentes}

\begin{table}[H]
  \centering
  \small
  \begin{tabularx}{\linewidth}{@{}>{\raggedright\arraybackslash}p{5.4cm}X@{}}
    \toprule
    \textbf{Síntoma} & \textbf{Solución} \\
    \midrule
    \emph{password authentication failed for user "postgres"}
      & Usar la contraseña definida al instalar PostgreSQL (paso A.1). \\
    \emph{psql no se reconoce como comando}
      & Usar la ruta completa de \bd{psql.exe}, como en el paso A.5. \\
    \emph{No such file or directory} al ejecutar \bd{psql} o \bd{docker cp}
      & PowerShell no está en la carpeta de los scripts: abrirlo desde esa
        carpeta (paso A.5) o escribir la ruta completa de cada archivo entre
        comillas. \\
    El puerto 5432 está ocupado durante la instalación
      & Aceptar el puerto que proponga el instalador; pgAdmin lo usa solo. En
        \bd{psql}, agregar \bd{-p} y el número de puerto. \\
    pgAdmin muestra una sola tabla de resultados al validar
      & Es el comportamiento normal con un archivo completo: ejecutar cada
        consulta por separado, como indica el paso A.4. \\
    En la opción B, \emph{port is already allocated}
      & El puerto 5432 ya lo usa otro PostgreSQL, por ejemplo el de la
        opción A. Ejecutar \bd{docker rm -f rapiditos-postgres} y repetir el
        paso B.1 con \bd{-p 127.0.0.1:5433:5432}. \\
    En la opción B, \emph{the database system is starting up} o
    \emph{is shutting down}
      & El servidor todavía se está iniciando: esperar el mensaje indicado en
        el paso B.1 y repetir el comando. \\
    En la opción B, \emph{The container name "/rapiditos-postgres" is already
    in use}
      & Ya existe un contenedor de una prueba anterior: eliminarlo con
        \bd{docker rm -f rapiditos-postgres} y repetir el paso B.1. \\
    \bd{docker} no responde o no encuentra el motor
      & Abrir Docker Desktop y esperar a que indique que está en ejecución. \\
    \bottomrule
  \end{tabularx}
\end{table}

\end{document}
